\documentclass[10pt,a4paper]{amsart}
\usepackage[margin=19mm]{geometry}
\usepackage{amsmath,amssymb,mathtools}
\usepackage[scr=rsfs,cal=euler]{mathalfa}
\usepackage{xfrac}
\usepackage[unicode,hidelinks,bookmarks=false]{hyperref}
\newcommand{\ie}{i.e.\ }
\newcommand{\cf}{cf.\ }
\newcommand{\resp}{respectively\ }
\newcommand{\Subsection}{\subsection}
\providecommand{\nobreakdash}{\nobreak\hbox{-}}
\setlength{\emergencystretch}{2em}
\multlinegap0pt
\usepackage{xcolor}
\usepackage{amssymb}
\usepackage[normalem]{ulem}
\usepackage{tikz}
\newtheorem{proposition}{Proposition}
\title{$C_4$-face-magic labeling on a $4 \times 4$ Klein \\ bottle grid graph}
\author{Timothy Myers, Stephen J. Curran}
\date{Source: arXiv:2606.06817v1 (CC BY 4.0)}
\begin{document}
\maketitle

\noindent\emph{Source and licence.} $C_4$-face-magic labeling on a $4 \times 4$ Klein \\ bottle grid graph. Version arXiv:2606.06817v1 (CC BY 4.0), distributed by arXiv under the Creative Commons Attribution 4.0 International licence, http://creativecommons.org/licenses/by/4.0/, as recorded in the arXiv licence metadata for this exact version and shown on the abstract page https://arxiv.org/abs/2606.06817v1. Full licence notice: \texttt{licenses/arxiv-2606.06817-CC-BY-4.0.txt}.\par\smallskip

\noindent\textit{Editorial scope.} Counted unit: the article's proposition prop:PermissibleA1Values (source0.tex lines 1160--1174). The article's complete original proof of it is delivered with it (source0.tex lines 1176--1458, one \texttt{proof} environment of 283 lines). No mathematics is written, repaired or re-derived: the statement and the proof are the article's own lines, in the article's own notation, and no retained line is substituted or re-flowed.  Retained uncounted as the source-local context the proof needs: lines 773--809.  Nothing was omitted from any retained span, so the package carries no omission record.  No retained line contains a citation, so the wrapper carries no bibliography and no bibliography entry.  No retained line contains a figure environment, an image-inclusion command or drawing code, so no image is delivered and the package declares no asset dependency.  Editorial formatting only, in addition: an AMS \texttt{amsart} preamble that restates the article's own declarations and macros used by the retained lines, namely the environments \texttt{proposition} on separate counters, together with the standard packages those lines need.  The retained environments print in this document as Proposition 1 in order of appearance, whereas the article prints the same items with the numbering of its own full document.  The complete native source of the article as distributed in the arXiv e-print for this exact version is bundled unchanged as \texttt{sources/202215/source0.tex}.
\begin{proposition}\label{prop:extendedgesums}
Suppose  $X=\{x_{i,j}: 1\le i \le 4 \text{ and } 
   1\le j \le 4\}$ is a $C_4$-face-magic labeling
   on $\mathcal{K}_{4,4}$ with face-magic value $S$.
   Let $a_1= x_{1,1} + x_{1,2}$,
    $a_2= x_{1,2} + x_{1,3}$, and
     $a_3= x_{1,4} + x_{1,1}$.
        Then
\begin{enumerate}
\item[(i)]
\begin{align*}
   a_1 &= x_{2i-1,2j-1} + x_{2i-1,2j} 
    \text{ \ \  and } \\
  S - a_1 &=  x_{2i,2j-1} + x_{2i,2j} 
    \text{ \ \ \ \ \ \ \ \ for } 1\le i \le 2 
    \text{ and } 1 \le j \le 2.
\end{align*}
\item[(ii)]
\begin{align*}
    a_2 & = x_{2i-1, 2} + x_{2i-1, 3}
    \text{ \ \  and } \\
    S - a_2 & = x_{2i, 2} + x_{2i,3}
    \text{ \ \ \ \ \ \ \ \ for } 1\le i \le 2
\end{align*}
\item[(iii)]
\begin{align*}
    a_3 & = x_{2i-1, 4} + x_{2i-1, 1}
    \text{ \ \  and } \\
    S - a_3 & = x_{2i, 4} + x_{2i,1}
    \text{ \ \ \ \ \ \ \ \ for } 1\le i \le 2
\end{align*}
\end{enumerate}
Furthermore, $2a_1 = a_2 + a_3$. In addition, for $1\le i \le 4$,
\begin{equation} \label{eqn:SameEdgeSum}
    x_{i,1} + x_{i,2} = x_{i,3} + x_{i,4}.
\end{equation}
\end{proposition}

\begin{proposition} \label{prop:PermissibleA1Values}
     Suppose $X=\{ x_{i,j} : (i,j) \in V(\mathcal{K}_{4,4})\}$ 
  is a properly symmetrized $C_4$-face-magic labeling on $\mathcal{K}_{4,4}$.
  Then the only permissible values of $a_1$ 
  are $a_1\in\{9,13,15,16,17\}$.
  Furthermore, for each permissible value of $a_1$,
  the corresponding permissible values of $a_2$ are given below.
  \begin{itemize}
    \item If $a_1=9$, then $a_2\in\{10,11,13\}$.
    \item If $a_1=13$, then $a_2\in\{14,15,21\}$.
    \item If $a_1=15$, then $a_2\in\{16,19,23\}$.
    \item If $a_1=16$, then $a_2\in\{18,20,24\}$.
    \item If $a_1=17$, then $a_2\in\{18,19,21,25\}$.
  \end{itemize}
\end{proposition}

\begin{proof}
     By Proposition \ref{prop:extendedgesums},
    the collection
\begin{equation}\label{partition1to16}
    \big\{ \{x_{2i-1,2j-1}, x_{2i-1,2j}\} : 1\le i,j\le 2\big\}
    \cup
     \big\{ \{x_{2i,2j-1}, x_{2i,2j}\} : 1\le i,j\le 2\big\}
\end{equation}
is a partition of the set $\{1,2,\ldots,16\}$ such that
$x_{2i-1,2j-1} + x_{2i-1,2j} = a_1$ and
$x_{2i,2j-1} + x_{2i,2j} =S- a_1$.
Thus, we need to partition the set $\mathcal{C}=\{1,2,\ldots,16\}$ into eight 2-element sets of distinct elements such that
  \begin{enumerate}
    \item[(i)]  the sum of the elements of a set is $a_1$ for four of the 2-element sets, and
    \item[(ii)]  the sum of the elements of a set  is $S-a_1$ for the other four 2-element sets.
  \end{enumerate}
   Note that $a_1\ge 9$ since the smallest positive integer which can be written as a sum of 2 distinct positive integers in exactly four different ways is $9$. Recall that $x_{1,1}=1$. Since $a_1=x_{1,1}+x_{1,2}=1+x_{1,2}$\ and $x_{1,2}\le 16$, then $a_1\le 17$.
  Thus the only possible values for $a_1$ are $9,10,11,12,13,14,15,16$ and $17$.\\
  \hspace{.1in}\\
  Consider the proof that $a_1$ can equal $9$.
  Since $9=1+8=2+7=3+6=4+5$, we will form four distinct two element sets with the integers in each of these four sums: $\{1,8\}, \{2,7\}, \{3, 6\}, \{4,5\}$; which satisfies (i). When these pairs are removed from $\mathcal{C}$ the remaining integers are 
  $9,10,11,12,13,14, 15$ and $16$. There is precisely one way to form four pairs of these integers such that their sum is $34-9=25$, which is $\{9, 16\}, \{10, 15\}, \{11, 14 \}$ and $ \{12, 13\}$; which satisfies (ii). Thus, $a_1$ can equal $9$.\\
  \hspace{2in}\\
  Similarly, as demonstrated in Figure \ref{permissiblevaluesa1} we can partition $\mathcal{C}$ so that (i) and (ii) hold for $a_1=13, 15, 16, 17$.\\
  \hspace{2in}\\

\begin{figure}[hbt!]
    \centering
\renewcommand{\arraystretch}{1.6}
\begin{tabular}{|c|c|c|c|} \hline
$a_1$  &  \text{Pairs in $\mathcal{C}$ that sum to $a_1$} &  $S-a_1$ & \text{Pairs in $\mathcal{C}$ that sum to $S-a_1$}\\  \hline
 $9$  &  $\{1,8\}, \{ 2, 7 \}, \{ 3, 6 \}, \{ 4, 5  \}$  & $25$ & $\{9,16\}, \{10, 15\}, \{  11, 14\}, \{ 12, 13 \}$ \\ \hline
 $13$ & $\{1,12\}, \{ 2, 11 \}, \{ 3, 10\}, \{ 4, 9  \}$ & $21$  & $\{5,16\}, \{ 6, 15 \}, \{ 7, 14\}, \{ 8, 13  \}$\\  \hline
 $15$ &  $\{1,14\}, \{ 2, 13 \}, \{ 5, 10 \}, \{ 6, 9  \}$ &  $19$  &  $\{7,12\}, \{ 8, 11 \}, \{ 3, 16 \}, \{ 4, 15  \}$ \\  \hline
 $16$  &  $\{1,15\}, \{ 3, 13 \}, \{ 5, 11\}, \{ 7, 9  \}$  &  $18$  & $\{2,16\}, \{ 4, 14 \}, \{ 6, 12 \}, \{ 8, 10  \}$  \\   \hline
 $17$ & $\{1,16\}, \{ 2, 15 \}, \{ 3, 14 \}, \{ 4, 13  \}$ & $17$  &  $\{5,12\}, \{ 6, 11 \}, \{ 7, 10\}, \{ 8, 9  \}$ \\ \hline
\end{tabular}
\caption{The four pairs of elements in $\mathcal{C}$ that sum to $a_1$ and $S-a_1$.}
    \label{permissiblevaluesa1}
\end{figure}
 We can eliminate $10, 11, 12$, and $14$ as possible values for $a_1$ as follows.
 Note that $a_1\neq 10$ as follows. Since we would be required to set $S-a_1=24$ if $a_1=10$, there is no integer in $\mathcal{C}$ which differs from $5$ that can be added to $5$ to total $10$ or $24$ since only $5+5=10$ and $5+19=24$.\\
  \hspace{2in}\\
To show that $a_1\neq 11$ we will establish a contradiction by assuming $a_1=11$. The pairs of distinct elements in $\mathcal{C}$ that sum to $a_1=11$ are 
\begin{equation}
    \{1, 10\}, \{2, 9  \}, \{ 3, 8 \}, \{ 4, 7  \}, \text{and}\ \{5, 6\};   \label{pairsumsa1}
\end{equation}
and the pairs of elements in $\mathcal{C}$ that sum to $S-a_1=23$ are  
\begin{equation}
   \{7, 16\}, \{8 ,15\}, \{9, 14\}, \{10, 13\}, \text{and}\ \{11, 12\}. \label{pairsumsS-a1} 
\end{equation} 
By \eqref{partition1to16} we need to choose four pairs from \eqref{pairsumsa1} and four pairs from \eqref{pairsumsS-a1} which when united form a partition of $\mathcal{C}$. Since $\{1, 10\}$ is the only pair with $1$ and $\{10, 13\}$ is the only pair with $13$ we need to include each of these pairs in the partition of $\mathcal{C}$. However, $\{1,10\}\cap \{10,13\}=\{10\}$, which contradicts the requirement that subsets of a partition are disjoint. Therefore $a_1\neq 11$. By similar reasoning, $a_1\neq 12, 14$.\\

We will now determine the permissible values for the edge sum $a_2$ for $a_1=9$. 
Since $x_{1,2}=8$ and $x_{1,3}\ge 2$ then $a_2=x_{1,2}+x_{1,3}\ge 10$. 
Note that since $x_{1,3}+x_{1,4}=9$ and $x_{1,4}\ge 2$ then $x_{1,3}\le 7$. 
Thus, $a_2=x_{1,2}+x_{1,3}\le 15$; so $a_2$ can potentially be $10, 11, 12, 13, 14$, or $15$.\\ 
\hspace{2in}\\
Of the sets $\{1,8\}$
$\{2,7\}$, $\{3,6\}$ and $\{4,5\}$,
let $A_1=\{1,8\}$, choose $A_2$ to be the set   
such that $x_{1,3}\in A_2$, and let $A_3$ and $A_4$ be the remaining two sets.
By Proposition \ref{prop:extendedgesums},
there are labels $y\in A_3$ and $z\in A_4$
such that $y+z=a_2$.
We test this criterion for the 
possible values of the label $x_{1,3}$ from
the set $\{2,3,4,5,6,7\}$.
In Figure  \ref{establish_a2_values}, we observe that
this criterion is only satisfied for $x_{1,3}\in \{2,3,5\}$.
Thus, $a_2\in\{10,11,13\}$. 
In Figure  \ref{eliminate_a2_values}, we see that 
this criterion is not satisfied for $x_{1,3}\in \{4,6,7\}$.
Thus, $a_2\not\in\{12,14,15\}$. 

\begin{figure}[hbt!]
    \centering
\renewcommand{\arraystretch}{1.5}
\begin{tabular}{|c|c|c|c|c|c|} \hline
$A_1$ & $A_2$  & $\textcolor{red}{x_{1,2}}+\textcolor{red}{x_{1,3}}=a_2$ : & $A_3$ &  $A_4$ & $\textcolor{red}{y}+\textcolor{red}{z} = a_2\ :$\\ 
      &            &   $\textcolor{red}{x_{1,2}}\in A_1$\ ,\ $\textcolor{red}{x_{1,3}}\in A_2$                                      &      &    &
$\textcolor{red}{y}\in A_3$\ ,\ $\textcolor{red}{z}\in A_4$ \\  \hline
$\{1, \textcolor{red}{8}\}$ & $\{\textcolor{red}{2}, 7\}$ & $\textcolor{red}{8}+\textcolor{red}{2}=10$ & $\{ 3,\textcolor{red}{6}\}$ &  $\{ \textcolor{red}{4}, 5  \}$ &  $\textcolor{red}{6}+\textcolor{red}{4}=10$\\  \hline
$\{1,\textcolor{red}{8}\}$ & $\{\textcolor{red}{3}, 6\}$ & $\textcolor{red}{8}+\textcolor{red}{3}=11$ & $\{ 2,\textcolor{red}{7}\}$ &  $\{ \textcolor{red}{4}, 5  \}$ &  $\textcolor{red}{7}+\textcolor{red}{4}=11$\\  \hline
$\{1, \textcolor{red}{8}\}$ & $\{4, \textcolor{red}{5}\}$ & $\textcolor{red}{8}+\textcolor{red}{5}=13$ & $\{ 2,\textcolor{red}{7}\}$ &  $\{ 3, \textcolor{red}{6}  \}$ &  $\textcolor{red}{7}+\textcolor{red}{6}=13$\\  \hline
\end{tabular}
    
    \caption{Using a criterion to establish $10, 11$, and $13$ as possible values for $a_2$ when $a_1=9$.}
    \label{establish_a2_values}
\end{figure}

\begin{figure}[hbt!]
\centering
\renewcommand{\arraystretch}{1.5}
\begin{tabular}{|c|c|c|c|c|c|} \hline
$A_1$ & $A_2$  & $\textcolor{red}{x_{1,2}}+\textcolor{red}{x_{1,3}}=a_2$ & $A_3$ &  $A_4$ & $y+z\neq a_2\ :$\\ 
      &            &   $\textcolor{red}{x_{1,2}}\in A_1$\ ,\ $\textcolor{red}{x_{1,3}}\in A_2$                                     &      &    &
$y\in A_3$\ ,\ $z\in A_4$ \\  \hline     
$\{1, \textcolor{red}{8}\}$ & $\{2,\textcolor{red}{7}\}$ & $\textcolor{red}{8}+\textcolor{red}{7}=15$ & $\{ 3,6\}$ &  $\{ 4, 5  \}$ &  $3+4\neq 15$\\ 
          &     &          &            &                 & $3+5\neq 15$\\ 
          &     &          &            &                 & $6+4\neq 15$\\   &     &          &            &                 & $6+5\neq 15$\\  \hline     
 $\{1, \textcolor{red}{8}\}$& $\{3,\textcolor{red}{6}\}$ & $\textcolor{red}{8}+\textcolor{red}{6}=14$ & $\{ 2,7\}$ &  $\{ 4, 5  \}$ &  $2+4\neq 14$\\ 
          &     &          &            &                 & $2+5\neq 14$\\ 
          &     &          &            &                 & $7+4\neq 14$\\   &     &          &            &                 & $7+5=14$\\     \hline
$\{1, \textcolor{red}{8}\}$ & $\{\textcolor{red}{4},5\}$ & $\textcolor{red}{8}+\textcolor{red}{4}=12$ & $\{ 2,7\}$ &  $\{ 3, 6  \}$ &  $2+3\neq 12$\\ 
          &     &          &            &                 & $2+6\neq 12$\\ 
          &     &          &            &                 & $7+3\neq 12$\\   &     &          &            &                 & $7+6\neq 12$\\     \hline          
\end{tabular}
\caption{Using a criterion to eliminate $15, 14$ and $12$ as possible values for $a_2$ when $a_1=9$ by summing all possible values for $y$ and $z$.}
\label{eliminate_a2_values}    
\end{figure}

The proofs for the permissible values of $a_2$ when $a_1$ is $13, 15$, and $16$ are similar to the case where $a_1=9$ (see the tables in Figures \ref{establish_a2_valuesais13}, \ref{establish_a2_valuesais15}, and  \ref{establish_a2_valuesais16}).

\begin{figure}[hbt!]
    \centering
\renewcommand{\arraystretch}{1.5}
\begin{tabular}{|c|c|c|c|c|c|} \hline
$A_1$ & $A_2$  & $\textcolor{red}{x_{1,2}}+\textcolor{red}{x_{1,3}}=a_2$ : & $A_3$ &  $A_4$ & $\textcolor{red}{y}+\textcolor{red}{z} = a_2\ :$\\ 
      &            &   $\textcolor{red}{x_{1,2}}\in A_1$\ ,\ $\textcolor{red}{x_{1,3}}\in A_2$                                      &      &    &
$\textcolor{red}{y}\in A_3$\ ,\ $\textcolor{red}{z}\in A_4$ \\  \hline
$\{1, \textcolor{red}{12}\}$ & $\{\textcolor{red}{2}, 11\}$ & $\textcolor{red}{12}+\textcolor{red}{2}=14$ & $\{ 3,\textcolor{red}{10}\}$ &  $\{ \textcolor{red}{4}, 9  \}$ &  $\textcolor{red}{10}+\textcolor{red}{4}=14$\\  \hline
$\{1,\textcolor{red}{12}\}$ & $\{\textcolor{red}{3}, 10\}$ & $\textcolor{red}{12}+\textcolor{red}{3}=15$ & $\{ 2,\textcolor{red}{11}\}$ &  $\{ \textcolor{red}{4}, 9  \}$ &  $\textcolor{red}{11}+\textcolor{red}{4}=15$\\  \hline
$\{1, \textcolor{red}{12}\}$ & $\{4, \textcolor{red}{9}\}$ & $\textcolor{red}{12}+\textcolor{red}{9}=21$ & $\{ 3,\textcolor{red}{10}\}$ &  $\{ 2, \textcolor{red}{11}  \}$ &  $\textcolor{red}{10}+\textcolor{red}{11}=21$\\  \hline
\end{tabular}
    
    \caption{Using a criterion to establish $10, 11$, and $13$ as possible values for $a_2$ when $a_1=13$.}
    \label{establish_a2_valuesais13}
\end{figure}

\begin{figure}[hbt!]
    \centering
\renewcommand{\arraystretch}{1.5}
\begin{tabular}{|c|c|c|c|c|c|} \hline
$A_1$ & $A_2$  & $\textcolor{red}{x_{1,2}}+\textcolor{red}{x_{1,3}}=a_2$ : & $A_3$ &  $A_4$ & $\textcolor{red}{y}+\textcolor{red}{z} = a_2\ :$\\ 
      &            &   $\textcolor{red}{x_{1,2}}\in A_1$\ ,\ $\textcolor{red}{x_{1,3}}\in A_2$                                      &      &    &
$\textcolor{red}{y}\in A_3$\ ,\ $\textcolor{red}{z}\in A_4$ \\  \hline
$\{1, \textcolor{red}{14}\}$ & $\{\textcolor{red}{2}, 13\}$ & $\textcolor{red}{14}+\textcolor{red}{2}=16$ & $\{ 5,\textcolor{red}{10}\}$ &  $\{ \textcolor{red}{6}, 9  \}$ &  $\textcolor{red}{10}+\textcolor{red}{6}=16$\\  \hline
$\{1,\textcolor{red}{14}\}$ & $\{\textcolor{red}{5}, 10\}$ & $\textcolor{red}{14}+\textcolor{red}{5}=19$ & $\{ \textcolor{red}{6}, 9 \}$ &  $\{2, \textcolor{red}{13}  \}$ &  $\textcolor{red}{6}+\textcolor{red}{13}=19$\\  \hline
$\{1, \textcolor{red}{14}\}$ & $\{6, \textcolor{red}{9}\}$ & $\textcolor{red}{14}+\textcolor{red}{9}=23$ & $\{ 2,\textcolor{red}{13}\}$ &  $\{ 5, \textcolor{red}{10}  \}$ &  $\textcolor{red}{13}+\textcolor{red}{10}=23$\\  \hline
\end{tabular}
    
    \caption{Using a criterion to establish $16, 19$, and $23$ as possible values for $a_2$ when $a_1=15$.}
    \label{establish_a2_valuesais15}
\end{figure}

\begin{figure}[hbt!]
    \centering
\renewcommand{\arraystretch}{1.5}
\begin{tabular}{|c|c|c|c|c|c|} \hline
$A_1$ & $A_2$  & $\textcolor{red}{x_{1,2}}+\textcolor{red}{x_{1,3}}=a_2$ : & $A_3$ &  $A_4$ & $\textcolor{red}{y}+\textcolor{red}{z} = a_2\ :$\\ 
      &            &   $\textcolor{red}{x_{1,2}}\in A_1$\ ,\ $\textcolor{red}{x_{1,3}}\in A_2$                                      &      &    &
$\textcolor{red}{y}\in A_3$\ ,\ $\textcolor{red}{z}\in A_4$ \\  \hline
$\{1, \textcolor{red}{15}\}$ & $\{\textcolor{red}{3}, 13\}$ & $\textcolor{red}{15}+\textcolor{red}{3}=18$ & $\{ 5,\textcolor{red}{11}\}$ &  $\{ \textcolor{red}{7}, 9  \}$ &  $\textcolor{red}{11}+\textcolor{red}{7}=18$\\  \hline
$\{1,\textcolor{red}{15}\}$ & $\{\textcolor{red}{5}, 11\}$ & $\textcolor{red}{15}+\textcolor{red}{5}=20$ & $\{ 3,\textcolor{red}{13}\}$ &  $\{ \textcolor{red}{7}, 9  \}$ &  $\textcolor{red}{13}+\textcolor{red}{7}=20$\\  \hline
$\{1, \textcolor{red}{15}\}$ & $\{7, \textcolor{red}{9}\}$ & $\textcolor{red}{15}+\textcolor{red}{9}=24$ & $\{ 5,\textcolor{red}{11}\}$ &  $\{ 3, \textcolor{red}{13}  \}$ &  $\textcolor{red}{11}+\textcolor{red}{13}=24$\\  \hline
\end{tabular}
    
    \caption{Using a criterion to establish $18, 20$, and $24$ as possible values for $a_2$ when $a_1=16$.}
    \label{establish_a2_valuesais16}
\end{figure}
  
  Suppose $a_1=17$. 
  Then $a_1=\tfrac{1}{2}S$ and $S-a_1=\tfrac{1}{2}S$.
  Thus, $X$ is vertically pairwise balanced.
  So, the sets
   $\{ x_{2i-1,2j-1},\, x_{2i-1,2j}\}$ and $\{ x_{2i,2j-1},\, x_{2i,2j}\}$,
 for $i,j=1,2$, are the eight sets $\{k,17-k\}$, for $1\le k \le 8$,
 which we denote by $A_1, \, A_2,\ldots,\, A_8$.
 By Proposition \ref{prop:extendedgesums}, $a_3=2a_1-a_2=S-a_2$. Then $S-a_3=a_2$.
 Hence,
   \begin{equation}\label{eqn:EdgeSumsAreA2}
      a_2=x_{2i-1,2}+x_{2i-1,3}=x_{2i,1}+x_{2i,4}
      \text{ \ for \ }i=1,\, 2.
  \end{equation}
 Thus, we can combine the sets $A_1, \, A_2,\ldots,\, A_8$ into pair sets
 $A_{2k-1},\, A_{2k}$, for $1\le k \le 4$, such that
there exist $y\in A_{2k-1}$ and $z\in A_{2k}$
  for which $y+z=a_2$.
  Let $A_1=\{ 1,16\}$.
  Then $A_2$ is the set that contains the element 
  $x\in A_2$ for which $x+16=a_2$.

Since $x_{1,2}=16$ and $x_{1,3}\ge 2$, we have
$a_2-x_{1,2}+x_{1,3}\ge 18$.
  By \eqref{eqn:EdgeSumsAreA2}, we have
  \begin{equation*}
      4a_2= \sum_{i=1}^2 (x_{2i-1,2}+x_{2i-1,3} +x_{2i,1}+x_{2i,4})
 \le \sum_{k=9}^{16} k =100.
  \end{equation*}
 Hence, $a_2\le 25$.

 Consider $a_2=18$. 
 Since $x_{1,1}=1$, we have $x_{1,2}=16$.
 Then $x_{1,3}=a_2-x_{1,2}=2$ and $A_2=\{2,\, 15\}$.
 Let $A_3=\{3,\, 14\}$. Since $3$ cannot match with $15$ to add to $a_2=18$,
 $14$ must match with $4$. Thus, $A_4=\{4, 13\}$.
 Let $A_5=\{5,\, 12\}$. Since $5$ cannot match with $13$ to add to $a_2=18$,
 $12$ must match with $6$. Thus, $A_6=\{6, 11\}$.
 Let $A_7=\{7,\, 10\}$. Since $7$ cannot match with $11$ to add to $a_2=18$,
 $10$ must match with $8$. Thus, $A_8=\{8, 6\}$.
 See line 2 of Figure \ref{fig:establish_a2_valuesFora1Is17}.
 A similar analysis establishes the sets $A_1,\, A_2,\ldots,\, A_8$
 for $a_2\in\{19,\, 21, \, 25\}$. See Figure \ref{fig:establish_a2_valuesFora1Is17}.
 
\begin{figure}[ht]
    \centering
    {\small
\renewcommand{\arraystretch}{1.5}
\begin{tabular}{|c|c|c|c|c|} \hline
$A_1, A_2$ & $\textcolor{blue}{x_{1,2}} +\textcolor{red}{x_{1,3}}=a_2$ & $A_3,\, A_4$ &  $A_5,\, A_6$ 
&  $A_7,\, A_8$\\  \hline
$\{1, \textcolor{blue}{16}\}, \{\textcolor{red}{2},15\}$ & $\textcolor{blue}{16}+\textcolor{red}{2}=18$ 
& $\{ 3,\textcolor{blue}{14}\}$, $\{ \textcolor{blue}{4},13\}$  
&  $\{ 5, \textcolor{blue}{12}  \}$, $\{ \textcolor{blue}{6}, 11\}$  
&  $\{7, \textcolor{blue}{10} \}$ , $\{\textcolor{blue}{8}, 9 \}$ \\  \hline
$\{1, \textcolor{blue}{16}\}, \{\textcolor{red}{3},14\}$ & $\textcolor{blue}{16}+\textcolor{red}{3}=19$ 
& $\{ 5,\textcolor{blue}{12}\}$, $\{ \textcolor{blue}{7},10\}$  
&  $\{ \textcolor{blue}{8}, 9  \}$, $\{6, \textcolor{blue}{11} \}$  
&  $\{\textcolor{blue}{4}, 13 \}$ , $\{2, \textcolor{blue}{15}\}$ \\  \hline
$\{1, \textcolor{blue}{16}\},   \{\textcolor{red}{5},12\}$ & $\textcolor{blue}{16}+\textcolor{red}{5}=21$ 
& $\{\textcolor{blue}{8},9\}$, $\{4, \textcolor{blue}{13}\}$  
&  $\{2, \textcolor{blue}{15}  \}$, $\{\textcolor{blue}{6}, 11 \}$  
&  $\{\textcolor{blue}{7},10 \}$ , $\{3, \textcolor{blue}{14}\}$ \\  \hline
$\{1, \textcolor{blue}{16}\},    \{8,\textcolor{red}{9}\}$ & $\textcolor{blue}{16}+\textcolor{red}{9}=25$ 
& $\{ 2,\textcolor{blue}{15}\}$, $\{7, \textcolor{blue}{10}\}$  
&  $\{3, \textcolor{blue}{14}\}$, $\{6, \textcolor{blue}{11} \}$  
&  $\{4, \textcolor{blue}{13} \}$ , $\{5, \textcolor{blue}{12}\}$ \\  \hline
\end{tabular}
    }
    \caption{Using a criterion to establish $18,\, 19, \, 21$, and $25$ as possible values for $a_2$.
    We have $A_1=\{1,\, \textcolor{blue}{16}\}$ where $\textcolor{blue}{x_{1,2}=16}$.
    The label $\textcolor{red}{x_{1,3}}\in A_2$ in highlighted in  \textcolor{red}{red}.
    For the pairs $A_{2k-1}, A_{2k}$, for $k=2,\, 3,\, 4$,
    the labels from each set that sum to $a_2$ are highlighted in 
    \textcolor{blue}{blue}.}
    
    \label{fig:establish_a2_valuesFora1Is17}
\end{figure}

Finally, we need to show $a_2\not\in\{ 20, 22, 23, 24\}$.
Suppose $a_2=20$. Then $x_{1,3}=a_2-x_{1,2}=6$ and $x_{1,4}=a_1-x_{1,3}=11$. Thus $A_1=\{ 1, 16\}$
and $A_2=\{ 4, 13\}$. Let $A_3=\{ 7,10\}$.
Since $10+10=20$, we cannot use use the label $10$ to
use to add to $a_2=20$. Since $7+13=20$, we have 
$A_4=\{ 4,13\}$. Hence, $A_2=\{ 4,13\}=A_4$ which contradicts $A_2\ne A_4$. See Figure \ref{fig:establish_a2_NonvaluesFora1Is17}.
The proofs of the cases $a_2\in\{ 22, 23, 24\}$
are similar. We leave the details to the reader.

\begin{figure}[ht]
    \centering
\renewcommand{\arraystretch}{1.5}
\begin{tabular}{|c|c|c|c|} \hline
$A_1, A_2$ & $\textcolor{blue}{x_{1,2}} +\textcolor{red}{x_{1,3}}=a_2$ & $A_3,\, A_4$ &  $A_5,\, A_6$ \\  \hline
$\{1, \textcolor{blue}{16}\}, \{\textcolor{red}{4},13\}$ & $\textcolor{blue}{16}+\textcolor{red}{4}=20$ 
&  $\{\textcolor{blue}{7}, 10 \}$ , 
$\{4, \textcolor{blue}{13} \}$ & \\  \hline


$\{1, \textcolor{blue}{16}\}, \{\textcolor{red}{6},11\}$ & $\textcolor{blue}{16}+\textcolor{red}{6}=22$ 
&  $\{3, \textcolor{blue}{14}  \}$, $\{\textcolor{blue}{8}, 9 \}$  
&  $\{4, \textcolor{blue}{13} \}$ , $\{8, \textcolor{blue}{9}\}$ \\  \hline


$\{1, \textcolor{blue}{16}\},   \{\textcolor{red}{7},10\}$ & $\textcolor{blue}{16}+\textcolor{red}{7}=23$ 
&  $\{4, \textcolor{blue}{13} \}$ , $\{7, \textcolor{blue}{10}\}$  & \\  \hline


$\{1, \textcolor{blue}{16}\},    \{\textcolor{red}{8}, 9\}$ & $\textcolor{blue}{16}+\textcolor{red}{8}=24$ 
& $\{ 2,\textcolor{blue}{15}\}$, $\{8, \textcolor{blue}{9}\}$ & \\  \hline
\end{tabular}
    
    \caption{Using a criterion to establish $20,\, 22, \, 23$, and $24$ are not possible values for $a_2$.
    We have $A_1=\{1,\, \textcolor{blue}{16}\}$ where $\textcolor{blue}{x_{1,2}=16}$.
    The label $\textcolor{red}{x_{1,3}}\in A_2$ in highlighted in  \textcolor{red}{red}.
    For the pairs $A_{2k-1}, A_{2k}$, for $k=2$ 
    (and $k=3$ if necessary),
    the labels from each set that sum to $a_2$ are highlighted in 
    \textcolor{blue}{blue}.
    Note that for each value of $a_2$, we are forced to duplicate a set $A_j$ which produces a contradiction.}
    \label{fig:establish_a2_NonvaluesFora1Is17}
\end{figure}
\end{proof}

\end{document}
